\chapter{Desarrollo del sistema}

\section{Arquitectura de software}

Se construyó un workspace ROS~2 Humble con los paquetes
\texttt{nova\_sm3\_description} y \texttt{nova\_gait\_controller}. El primero
contiene el modelo, configuraciones de controladores y lanzadores de Gazebo y
MuJoCo. El segundo reúne la cinemática, generación de marchas, métricas,
supervisor de seguridad, monitor de contactos, comparador de fase y monitor de
estabilidad.

El controlador recibió órdenes de alto nivel en
\texttt{/nova/gait\_command} y publicó referencias de doce articulaciones en
\texttt{/joint\_trajectory\_controller/joint\_trajectory}. Cada referencia se
acompañó con un mensaje de fase que identificó marcha, muestra, ciclo, pata en
oscilación, subfase y contactos previstos. Esta sincronización permitió
segmentar los ensayos y comparar el plan con la respuesta simulada.

La Figura~\ref{fig:arquitectura-software} resume el flujo implementado. Las
órdenes de alto nivel generaron referencias y fase; el simulador produjo pose,
estados articulares, IMU y contactos; los nodos de observación calcularon
métricas y diagnósticos que fueron registrados y puestos a disposición del
supervisor.

\begin{figure}[H]
\centering
\resizebox{0.98\textwidth}{!}{%
\begin{tikzpicture}[
node distance=1.1cm and 1.35cm,
block/.style={rectangle,draw,rounded corners,fill=blue!12,
text width=3.15cm,minimum height=1.1cm,align=center,font=\small},
data/.style={rectangle,draw,rounded corners,fill=green!12,
text width=3.15cm,minimum height=1.1cm,align=center,font=\small},
arrow/.style={-latex,semithick}]
\node[block] (cmd) {Orden de marcha};
\node[block,right=of cmd] (gait) {Controlador de\\marcha e IK};
\node[block,right=of gait] (sim) {Gazebo / MuJoCo};
\node[data,below=of gait] (phase) {Fase y contactos\\previstos};
\node[data,below=of sim] (sensors) {Pose, articulaciones,\\IMU y contactos};
\node[block,below=of phase] (analysis) {Métricas, estabilidad\\y comparación};
\node[block,left=of analysis] (safety) {Supervisor de seguridad};
\node[data,right=of analysis] (bag) {rosbag2, CSV,\\gráficas e informes};
\draw[arrow] (cmd) -- (gait);
\draw[arrow] (gait) -- (sim);
\draw[arrow] (gait) -- (phase);
\draw[arrow] (sim) -- (sensors);
\draw[arrow] (phase) -- (analysis);
\draw[arrow] (sensors) -- (analysis);
\draw[arrow] (analysis) -- (safety);
\draw[arrow] (analysis) -- (bag);
\draw[arrow] (safety) |- (cmd);
\end{tikzpicture}%
}
\caption{Arquitectura funcional implementada para generación, simulación,
monitoreo, seguridad y registro de evidencia. Fuente: elaboración propia.}
\label{fig:arquitectura-software}
\end{figure}

\section{Temporización reproducible}

La implementación inicial avanzaba cada fase desde el instante tardío del
callback. Como consecuencia, una duración configurada de 0,18~s se ejecutaba
cerca de 0,20~s y el ciclo de gateo alcanzaba aproximadamente 4,80~s. Se
corrigió el planificador para avanzar desde el vencimiento programado anterior.
Una validación de seis ciclos midió 0,180003~s por fase y 4,320107~s por ciclo.
Las bolsas anteriores se conservaron, pero se separaron de la nueva versión.

\section{Marcha de gateo}

El gateo cartesiano coordinó las patas en el orden FL--RR--FR--RL. La trayectoria
incluyó transferencia de peso, precarga, despegue, vuelo, aterrizaje y contacto
final. Se mantuvieron configurables la transferencia lateral y longitudinal y
las curvas de ascenso y descenso. Los parámetros nominales empleados en los
resultados fueron paso de 18~mm, elevación de 14~mm y ciclo de 4,32~s.

Durante el desarrollo se ensayaron curvas alternativas para corregir el contacto
trasero. Los candidatos se evaluaron primero por alcanzabilidad, periodicidad y
salto articular inferior a 0,20~rad. El aumento de altura temprana trasera hasta
0,80 produjo una mejora temporal pequeña y se rechazó; valores desde 0,85
excedieron el límite articular adoptado.

\section{Marcha paso}

Para evitar que ambos simuladores usaran parámetros nominalmente equivalentes
pero codificados por separado, se añadió
\texttt{digital\_twin\_profiles.yaml}. Un generador produce el URDF y mundo SDF
de Gazebo, el MJCF de MuJoCo y un manifiesto con hashes. Una auditoría rechaza
el par si masa total, amortiguamiento, fricción o límite de esfuerzo no
coinciden. Se generaron y auditaron los perfiles nominal, realista provisional
y masa/fricción reducidas. La comparación dinámica completa permanece abierta
porque la interfaz MuJoCo todavía no expone pose corporal y contactos con el
mismo contrato de Gazebo.

Se implementó una marcha paso independiente con 32 referencias por ciclo,
longitud de 16~mm, elevación de 8~mm y transferencia lateral de 4~mm. El orden
de las patas permaneció FL--RR--FR--RL, pero la transferencia de peso y la
trayectoria se parametrizaron por separado del gateo. La marcha se verificó en
Gazebo y, para cadencia y articulaciones, en MuJoCo.

\section{Monitoreo y seguridad}

El nodo de métricas publicó posición, altura, orientación y estados articulares.
El supervisor ordenó postura segura ante estímulos provocados de altura,
inclinación, contacto, margen, vencimiento de datos, límite articular y
discontinuidad. Un nodo de prueba usó tópicos aislados y conservó en JSON la
activación, el motivo y la orden \texttt{stand}. Las entradas de contacto,
vencimiento y margen se mantuvieron deshabilitadas por defecto mientras se
define el rearme seguro y se cuantifican falsos positivos.

\section{Preparación del computador embebido}

Se instaló Ubuntu 22.04.5 Desktop ARM64 en una Raspberry Pi~4 y se verificaron
Wi-Fi, SSH y comunicación DDS con \texttt{ROS\_DOMAIN\_ID=42}. El workspace se
compiló en la Raspberry y se detectó el PCA9685 en la dirección I$^2$C 0x40. Se
realizaron pruebas progresivas de canales con alimentación externa; estas
pruebas no constituyeron calibración articular ni autorizaron una marcha
física.

Dos servomotores fueron sustituidos y se reforzaron acoples mecánicos. Debido a
que faltaba completar e inspeccionar el ensamble, no se presentaron mediciones
geométricas parciales como caracterización definitiva.

\section{Contrato de corrección residual aprendida}

Se implementó una interfaz independiente del simulador para limitar una futura
política de aprendizaje por refuerzo. La acción consta de doce correcciones
articulares acotadas a $\pm0{,}08$ rad, con variación máxima de $0{,}02$ rad por
paso. La corrección se suma a la marcha nominal y vuelve a limitarse con la
envolvente articular; no produce PWM ni puede anular el supervisor. La función
de recompensa penaliza inclinación, error de altura, error articular, magnitud
de acción y activaciones de seguridad. Esta interfaz fue probada, pero aún no
es una política entrenada ni un entorno dinámico completo.

Como preparación del entrenamiento se implementó un banco PPO residual en
NumPy con las cinco semillas congeladas (11, 23, 37, 53 y 71). El entorno
reducido reproduce observaciones, saturaciones y terminaciones, pero no resuelve
contacto ni dinámica de Gazebo; por ello sus políticas y curvas de aprendizaje
son evidencia de integración preparatoria y no una mejora locomotora validada.
